\Article{16}
First, we have $\Sigma Tex=T\cos u\, .\, \Sigma ae-T\sin u\, .\, \Sigma be=mT\cos u$ if we denote the \SourceLeaf{107}moment of the free magnetism of the first needle $\Sigma ae$ by $m$, since we know that $\Sigma be=0$. The contribution to $\frac{\mathrm{d}\Omega}{\mathrm{d}u}$ from the first term of $\Omega$ will be $=-mT\sin u$.

If, for brevity, we set
\[
\begin{aligned}
k={}&\alpha\cos\psi+\beta\sin\psi+A\cos(\psi-U)+B\sin(\psi-U)\\
&-a\cos(\psi-u)-b\sin(\psi-u)\\
l={}&(\alpha\sin\psi-\beta\cos\psi+A\sin(\psi-U)-B\cos(\psi-U)\\
&-a\sin(\psi-u)+b\cos(\psi-u))^2+(C-c)^2
\end{aligned}
\]
then $rr=(R+k)^2+l$.

Since in useful experiments $R$ must be much larger than the dimensions of either needle, the quantity $\frac1{r^{n-1}}$ expands into the rapidly convergent series
\[
\begin{aligned}
&R^{-(n-1)}-(n-1)kR^{-n}+\left(\frac{nn-n}{2}kk-\frac{n-1}{2}l\right)R^{-(n+1)}\\
&\qquad-\left(\tfrac16(n^3-n)k^3-\tfrac12(nn-1)kl\right)R^{-(n+2)}+\text{etc.}
\end{aligned}
\]
The rule for this series could easily be given if it were worth the trouble. Each term of the sum $\Sigma\frac{eE}{r^{n-1}}$ that results from substituting the values of $k,l$ will contain a factor of the form
\[
\Sigma eEa^\lambda b^\mu c^\nu A^{\lambda'}B^{\mu'}C^{\nu'}
\]
This factor equals the product of the factors $\Sigma ea^\lambda b^\mu c^\nu$, $\Sigma EA^{\lambda'}B^{\mu'}C^{\nu'}$, which depend respectively on the magnetic states of the first and second needles. The general conclusions we can draw in this regard are limited to the equations
\[
\begin{gathered}
\Sigma e=0,\ \Sigma ea=m,\ \Sigma eb=0,\ \Sigma ec=0,\\
\Sigma E=0,\ \Sigma EA=M,\ \Sigma EB=0,\ \Sigma EC=0
\end{gathered}
\]
where we denote by $M$ the moment of the free magnetism of the second needle. Consider the special case in which the shape of the first needle and the distribution of its magnetism are symmetric along its length, so that elements always correspond in pairs with opposite values of $a$ and $e$ and equal values of $b$ and $c$, with the center coinciding with the point $h$. Then $\Sigma ea^\lambda b^\mu c^\nu=0$ for every even value of the number $\lambda+\mu+\nu$. The same applies to the second needle if its shape and the distribution of its magnetism are symmetric about the point $H$. In general, therefore, the coefficients in the sum $\Sigma\frac{eE}{r^{(n-1)}}$ of the powers $R^{-(n-1)}$ and $R^{-n}$ vanish. If both needles are symmetric and symmetrically magnetized, and if the center of the first coincides with both $h$ and $h'$ and the center of the second with $H$, then the coefficients of the powers $R^{-(n+2)},R^{-(n+4)},R^{-(n+6)}$ and so on also vanish. Whenever those conditions hold \SourceLeaf{108}approximately, these coefficients must at least be very small. The leading term arising from the expansion of the second part of $\Omega$, namely $-\Sigma\frac{eE}{(n-1)r^{(n-1)}}$, will be
\[
\begin{aligned}
&=-\tfrac12 R^{-(n+1)}(n\Sigma eEkk-\Sigma eEl)\\
&=mMR^{-(n+1)}(n\cos(\psi-U)\cos(\psi-u)-\sin(\psi-U)\sin(\psi-u))
\end{aligned}
\]

It follows that the part of $\frac{\mathrm{d}\Omega}{\mathrm{d}u}$ corresponding to the action of the second needle is expressed by a series of the form
\[
fR^{-(n+1)}+f'R^{-(n+2)}+f''R^{-(n+3)}+\text{etc.}
\]
The coefficients are rational functions of the cosines and sines of the angles $\psi,u,U$ and of the quantities $\alpha,\beta$. They also involve constant quantities that depend on the magnetic state of the needles. In particular, we have
\[
f=mM(n\cos(\psi-U)\sin(\psi-u)+\sin(\psi-U)\cos(\psi-u))
\]
For our purpose we do not need the full expansion of the subsequent coefficients $f',f''$ and so on. It is enough to observe the following.

1) In the case of perfect symmetry just described, the coefficients $f',f'''$ and so on vanish.

2) If, with all other quantities unchanged, $\psi$ is increased by two right angles (or, equivalently, if the distance $R$ is taken on the same straight line extended backward on the other side of the point $h'$), the coefficients $f,f'',f''''$ and so on retain their values, while $f',f''',f^{\mathrm{v}}$ and so on take the opposite values. Thus the series becomes
\[
fR^{-(n+1)}-f'R^{-(n+2)}+f''R^{-(n+3)}-\text{etc.}
\]
This is readily seen from the fact that this change in $\psi$ takes $k$ to $-k$ while leaving $l$ unchanged.

